\section{复现说明与关键执行步骤}
计算使用Python、NumPy、SciPy及HiGHS，结果表通过OpenPyXL回读。
运行环境提供 uv 与 requirements.txt 两种安装方式（详见支撑材料 README.txt）；正式Q3无场景随机性，对冲扩展种子为7、蒙特卡洛种子为42。复现时先进入随附Python工程目录。
命令及结果对应关系见表\ref{tab:repro}，求解中的微小正则项不计入正文购电费用。
\begin{table}[htbp]\centering
\caption{核心复现入口}\label{tab:repro}
\small
\begin{tabularx}{\textwidth}{p{.60\textwidth}X}
\toprule
命令 & 对应内容\\\midrule
\texttt{uv run python -m solve.q1} & 典型日调度\\
\texttt{uv run python -m solve.q2\_tune --result2} & 问题二正式结果\\
\texttt{uv run python -m solve.q3} & 问题三正式结果\\
\texttt{uv run python -m solve.q4\_price\_fit} & 历史电价成本标定\\
\texttt{uv run python -m solve.q4 --result4-2} & 波动价格无调整层\\
\texttt{uv run python -m solve.q4 --result4-3} & 波动价格调整层\\
\texttt{uv run python -m solve.q3\_ablation} & 正式方案与对冲扩展消融\\
\bottomrule
\end{tabularx}
\end{table}
复现顺序应先形成历史成本表与权重，再执行问题二和问题三，最后重算电价标定与问题四。
同一随机种子有助于比较方案，但线性规划的退化解、求解器版本和浮点阈值仍可能造成小规模数值差异。
\begin{lstlisting}[language=Python,basicstyle=\ttfamily\footnotesize,breaklines=true]
for day in evaluation_days:
    weights = update_from_past_costs(day)
    forecast = predict_two_days(day, weights)
    plan = optimize_horizon(forecast, actual_soc)
    for slot in current_day:
        if forecast_is_released(slot):
            commitment = adjust_remaining_plan(
                released_forecast, actual_soc, plan)
        load, pv = observe_current_slot(slot)
        charge, discharge, emergency, spill = execute(
            commitment[slot], load, pv, actual_soc)
        actual_soc += eta * charge - discharge / eta
    settle_day_and_export()
\end{lstlisting}
上述为算法逻辑说明，完整求解、时间映射和回读审计以随附源程序为准。
论文指定日期表直接取自正式Excel，计划量与调整量分别保留；紧急购电表列出全部对应事件，无事件日明确记为零。

\section{AI辅助说明}
本次论文整理使用AI编码助手辅助核对已有结果、组织公式和段落、生成表格与LaTeX排版。
数值依据为给定附件及随附求解程序输出，文献条目经公开来源核对。
此说明记录本次辅助范围，不替代参赛者对题意、模型与最终提交材料的审阅。
